% !PRINTED 700
\weberpage{700}
%<<CONT
1. \gesperrt{Der Polygonquotient (1) ist ein Differential,
wenn die Klassen $A$, $B$ von $\mathfrak{A}$ und $\mathfrak{B}$ relativ prim
sind und der Bedingung genügen:}
\weq{2}{A = WB.}

Sind $a$ und $b$ die Ordnungszahlen von $A$ und $B$, so ist:
\weq{3}{a - b = 2p - 2,}
und da $A$ jedenfalls von der zweiten Gattung ist, nach \S\,199, (15):
\weq{4}{(O, A) = b + p - 1.}

Besteht $B$ aus einem einzelnen Punkt $\mathfrak{P}$, so ist $(O, A) = p$,
und die Dimension von $A$ ist ebenso gro\ss\ wie die von $W$; ist daher
\[
(\mathfrak{W}_1, \mathfrak{W}_2, \ldots, \mathfrak{W}_p)
\]
eine Basis von $W$, so ist
\[
(\mathfrak{P} \mathfrak{W}_1, \mathfrak{P} \mathfrak{W}_2, \ldots, \mathfrak{P} \mathfrak{W}_p)
\]
eine Basis von $A$. Daher ist $A$ eine uneigentliche Klasse mit
dem Teiler $\mathfrak{P}$, und es folgt:

2. \gesperrt{Ein einzelner Punkt $\mathfrak{P}$ kann nicht Untereck
eines Differentials sein.}

Hat die Klasse $A$ einen Teiler $\mathfrak{M}$ vom Grade $m$, so mu\ss\ $\mathfrak{M}$
im Teiler von $B$ enthalten sein; denn $\mathfrak{M}$ mu\ss\ in jedem der
Polygone $\mathfrak{W}\mathfrak{B}$, und folglich, da $W$ eine eigentliche Klasse ist, in
jedem der $\mathfrak{B}$ aufgehen. Ist
\[
\mathfrak{M} A' = A, \quad \mathfrak{M} B' = B,
\]
so mu\ss
\weq{5}{(O, A) = (O, A')}
sein. Es ist aber nach (3) und (4), auf $A$ und $A'$ angewandt,
wenn $A'$ von der zweiten Gattung ist:
\[
(O, A) = (O, A') + m,
\]
also $(O, A) > (O, A')$; folglich mu\ss\ $m = 0$ sein.

$A'$ kann aber nur dann von der ersten Gattung sein, wenn
$B$ ein isoliertes Polygon und $B = \mathfrak{M}$ ist. Dann ist:
\[
A' = W, \quad (O, A') = p,
\]
also mu\ss\ $(O, A) = p$, $b = 1$ sein, und wir kommen auf den
Fall des Satzes 2. Folglich:

3. \gesperrt{Jedes Polygon von mehr als einem Punkt kann
Untereck eines Differentials sein.}

Unter einem \gesperrt{Differential zweiter Gattung} versteht man
ein solches, dessen Untereck eine Potenz eines einzelnen Punktes $\mathfrak{P}$
%CONT>>
